% =============================================================================
% UBS COE EXPERIENCE DESIGN - FRAGE-ABO OFFERTE (STANDALONE)
% =============================================================================
% FehrAdvice & Partners AG
% Compile with: pdflatex UBS_037_frage_abo_coe_xd_proposal_standalone.tex
% =============================================================================

\documentclass[11pt,a4paper]{article}

% =============================================================================
% PACKAGES
% =============================================================================
\usepackage[T1]{fontenc}
\usepackage[utf8]{inputenc}
\usepackage[ngerman]{babel}
\usepackage[top=25mm,bottom=20mm,left=25mm,right=20mm]{geometry}
\usepackage{xcolor}
\usepackage{helvet}
\usepackage{tcolorbox}
\usepackage{booktabs}
\usepackage{colortbl}
\usepackage{hyperref}
\usepackage{titlesec}
\usepackage{fancyhdr}
\usepackage{lastpage}
\usepackage{graphicx}
\usepackage{enumitem}
\usepackage{array}
\usepackage{amssymb}

% =============================================================================
% COLORS (FehrAdvice Corporate Identity)
% =============================================================================
\definecolor{fadarkblue}{RGB}{2,64,121}
\definecolor{falightblue}{RGB}{84,158,222}
\definecolor{fadarkgray}{RGB}{37,33,42}
\definecolor{falightgray}{RGB}{243,245,247}
\definecolor{falilac}{RGB}{161,160,198}
\definecolor{famint}{RGB}{126,189,172}
\definecolor{faocher}{RGB}{222,203,63}
\definecolor{faorange}{RGB}{222,157,62}
\definecolor{fawhite}{RGB}{255,255,255}

% =============================================================================
% FONT SETUP (Helvetica fallback)
% =============================================================================
\renewcommand{\familydefault}{\sfdefault}
\AtBeginDocument{\color{fadarkgray}}

% =============================================================================
% HEADING STYLES
% =============================================================================
\titleformat{\section}{\bfseries\Large\color{fadarkblue}}{\thesection}{1em}{}
\titleformat{\subsection}{\bfseries\large\color{fadarkblue}}{\thesubsection}{1em}{}
\titleformat{\subsubsection}{\bfseries\normalsize\color{fadarkblue}}{\thesubsubsection}{1em}{}

% =============================================================================
% HYPERREF
% =============================================================================
\hypersetup{
    colorlinks=true,
    linkcolor=fadarkblue,
    citecolor=fadarkblue,
    urlcolor=fadarkblue,
    pdftitle={Verhaltensökonomisches Frage-Abo für CoE Experience Design},
    pdfauthor={FehrAdvice \& Partners AG}
}

% =============================================================================
% HEADER & FOOTER
% =============================================================================
\pagestyle{fancy}
\fancyhf{}
\fancyhead[L]{\small\color{fadarkgray}Verhaltensökonomisches Frage-Abo | CoE Experience Design}
\fancyhead[R]{\small\color{fadarkgray}Januar 2026}
\fancyfoot[L]{\small\color{fadarkblue}FehrAdvice \& Partners AG}
\fancyfoot[R]{\small\color{fadarkgray}Seite \thepage\ von \pageref{LastPage}}
\renewcommand{\headrulewidth}{0.5pt}
\renewcommand{\footrulewidth}{0pt}

\fancypagestyle{plain}{
    \fancyhf{}
    \fancyfoot[L]{\small\color{fadarkblue}FehrAdvice \& Partners AG}
    \fancyfoot[R]{\small\color{fadarkgray}Seite \thepage\ von \pageref{LastPage}}
    \renewcommand{\headrulewidth}{0pt}
}

% =============================================================================
% BOXES
% =============================================================================
\newtcolorbox{fainfobox}[1][]{
    colback=falightgray,
    colframe=fadarkblue,
    coltitle=fawhite,
    fonttitle=\bfseries,
    title=#1
}

\newtcolorbox{fawarningbox}[1][Hinweis]{
    colback=faorange!10,
    colframe=faorange,
    coltitle=fawhite,
    fonttitle=\bfseries,
    title=#1
}

\newtcolorbox{fasuccessbox}[1][]{
    colback=famint!10,
    colframe=famint,
    coltitle=fawhite,
    fonttitle=\bfseries,
    title=#1
}

% =============================================================================
% TABLE HELPERS
% =============================================================================
\newcommand{\fahead}[1]{\textcolor{fawhite}{\bfseries #1}}

% =============================================================================
% DOCUMENT
% =============================================================================

\begin{document}

% =============================================================================
% TITLE PAGE
% =============================================================================
\begin{titlepage}
\begin{center}
\vspace*{2cm}

% Logo placeholder
{\Huge\color{fadarkblue}\textbf{FehrAdvice}}\\[0.3cm]
{\large\color{falightblue}Behavioral Economics}

\vspace{2cm}

% Title
{\Huge\bfseries\color{fadarkblue}Verhaltensökonomisches\\[0.3cm]Frage-Abo\par}

\vspace{1cm}

% Subtitle
{\Large\color{falightblue}für das Center of Excellence Experience Design\par}

\vspace{1cm}

% Line
{\color{fadarkblue}\rule{\textwidth}{2pt}}

\vspace{2cm}

% Client
{\large UBS Switzerland AG\par}

\vspace{0.5cm}

% Date
{\large Januar 2026\par}

\vfill

% Footer
{\small\color{fadarkgray}
FehrAdvice \& Partners AG\\
Klausstrasse 4 | 8008 Zürich\\
www.fehradvice.com
}

\vspace{1cm}
{\small\color{fadarkgray}\textit{Vertraulich}}

\end{center}
\end{titlepage}

% =============================================================================
% TABLE OF CONTENTS
% =============================================================================
\tableofcontents
\newpage

% =============================================================================
\section{Executive Summary}
% =============================================================================

\subsection{Hintergrund und Problemstellung}

Das Center of Excellence for Experience Design (CoE XD) hat sich als zentrale Einheit für kundenzentrierte Lösungen bei UBS Switzerland etabliert. Mit dem erfolgreich abgeschlossenen Capability Building Projekt (UBS\_036) verfügt das Team nun über ein solides Fundament in verhaltensökonomischen Methoden, Storytelling und Experimentation.

Die Herausforderung besteht nun darin, dieses Wissen im Arbeitsalltag kontinuierlich anzuwenden und weiterzuentwickeln. Im täglichen Geschäft entstehen laufend Fragestellungen zu Customer Journeys, Touchpoint-Optimierungen und Design-Entscheidungen, die von einer verhaltensökonomischen Perspektive profitieren würden.

\subsection{Projektziel und Ansatz}

\begin{fainfobox}[Ziel des Angebots]
Etablierung eines strukturierten Prozesses für monatliche verhaltensökonomische Analysen, der:
\begin{itemize}
    \item Die im Capability Building erlernten Methoden in der Praxis verankert
    \item Dem CoE XD Team einen kontinuierlichen Sparringspartner bietet
    \item Design-Entscheidungen auf eine wissenschaftlich fundierte Basis stellt
    \item Die strategische Positionierung des CoE XD innerhalb von UBS stärkt
\end{itemize}
\end{fainfobox}

\subsection{Erwartete Ergebnisse und Mehrwerte}

\begin{itemize}
    \item \textbf{Für das CoE XD Team:} Schnelle, fundierte Antworten auf Experience-Design-Fragen -- von Journey-Optimierungen bis zu Touchpoint-Entscheidungen
    \item \textbf{Für UBS Kund:innen:} Bessere Kundenerlebnisse durch evidenzbasiertes Design
    \item \textbf{Für UBS Switzerland:} Stärkung des CoE XD als strategischer Partner mit messbarem Impact
\end{itemize}

\newpage
% =============================================================================
\section{Ausgangslage und Herausforderung}
% =============================================================================

\subsection{Aufbauend auf dem Capability Building}

Mit dem Projekt UBS\_036 hat das CoE XD wichtige Grundlagen geschaffen:

\begin{itemize}
    \item Ein gemeinsames Verständnis verhaltensökonomischer Prinzipien
    \item Erste BE Ambassadors, die als interne Multiplikatoren wirken
    \item Storytelling-Fähigkeiten zur Kommunikation von Design-Entscheidungen
    \item Ein Fundament für das langfristige Curriculum
\end{itemize}

\begin{fainfobox}[Die nächste Stufe]
Diese Fähigkeiten müssen nun im Alltag angewendet und vertieft werden. Dafür braucht es einen strukturierten Rahmen, der sicherstellt, dass verhaltensökonomische Analysen nicht ad hoc, sondern systematisch erfolgen.
\end{fainfobox}

\subsection{Warum ein strukturierter Analyseprozess notwendig ist}

Im Arbeitsalltag des CoE XD entstehen laufend Fragestellungen:

\begin{center}
\begin{tabular}{p{4cm}p{9cm}}
\toprule
\rowcolor{fadarkblue}
\fahead{Fragetyp} & \fahead{Beispiele} \\
\midrule
\rowcolor{fawhite}
Journey-Design & Wo in der Customer Journey verlieren wir Kund:innen? \\
\rowcolor{falightgray}
Touchpoint-Optimierung & Wie sollte ein Onboarding-Flow gestaltet sein? \\
\rowcolor{fawhite}
Choice Architecture & Wie viele Optionen sind zu viel? \\
\rowcolor{falightgray}
Behavioral Barriers & Warum nutzen Kund:innen Feature X nicht? \\
\rowcolor{fawhite}
A/B-Test Design & Welche Varianten sollten wir testen? \\
\bottomrule
\end{tabular}
\end{center}

\vspace{0.5cm}

Eine systematische verhaltensökonomische Analyse kann:

\begin{itemize}
    \item \textbf{Versteckte Barrieren aufdecken}, die in traditionellen UX-Analysen übersehen werden
    \item \textbf{Wirkungsprognosen liefern}, bevor Ressourcen in die Umsetzung fliessen
    \item \textbf{Design-Entscheidungen absichern} mit wissenschaftlicher Fundierung
    \item \textbf{Interne Akzeptanz stärken}, weil Empfehlungen nachvollziehbar begründet sind
\end{itemize}

\newpage
% =============================================================================
\section{Zielsetzung}
% =============================================================================

\subsection{Strategische Ziele}

\begin{itemize}
    \item[$\checkmark$] \textbf{Capability Building fortführen:} Die im Projekt UBS\_036 erlernten Methoden werden im Alltag verankert und kontinuierlich weiterentwickelt.
    \item[$\checkmark$] \textbf{Strategische Positionierung stärken:} Das CoE XD etabliert sich als Quelle evidenzbasierter Design-Empfehlungen.
    \item[$\checkmark$] \textbf{Messbare Wirkung nachweisen:} Durch dokumentierte Analysen wird der Wertbeitrag des CoE XD sichtbar.
\end{itemize}

\subsection{Verhaltensorientierte Ziele}

\begin{itemize}
    \item[$\checkmark$] \textbf{Wissenschaftlich fundierte Design-Entscheidungen:} Fragestellungen werden mit verhaltensökonomischen Modellen analysiert.
    \item[$\checkmark$] \textbf{Behavioral Barriers identifizieren:} Analysen zeigen versteckte Barrieren auf.
    \item[$\checkmark$] \textbf{Konsistenz durch systematische Methodik:} Analysen beruhen auf dem EBF Framework.
\end{itemize}

\subsection{Operative Ziele}

\begin{itemize}
    \item[$\checkmark$] \textbf{Pragmatische Analysen für den Alltag:} Jede Fragestellung führt zu konkreten Design-Empfehlungen.
    \item[$\checkmark$] \textbf{Frühwarnsystem für UX-Risiken:} Potenzielle Friction Points werden frühzeitig identifiziert.
    \item[$\checkmark$] \textbf{Kontinuierliches Lernen etablieren:} Wissensbasis wächst über Zeit.
\end{itemize}

\newpage
% =============================================================================
\section{Unser Vorschlag zum Vorgehen}
% =============================================================================

\subsection{Phase 1: Definition der Fragestellungen}

\begin{itemize}
    \item Das CoE XD definiert monatlich die relevanten inhaltlichen Fragestellungen
    \item Die Fragestellungen werden an FehrAdvice übermittelt mit allen notwendigen Unterlagen
    \item Festlegung: \textbf{Kurzanalyse} (3--5 Seiten) oder \textbf{umfassende Analyse} (8--10 Seiten)
\end{itemize}

\subsection{Phase 2: Verhaltensökonomische Analyse}

Systematische Analyse entlang verhaltensökonomischer Logik:

\begin{itemize}
    \item \textbf{Kontextfaktoren ($\Psi$-Dimensionen):} Situative Faktoren
    \item \textbf{Behavioral Barriers:} Kognitive, emotionale und praktische Hürden
    \item \textbf{Nutzenfunktionen:} Relevante Utility-Dimensionen (FEPSDE)
    \item \textbf{Intervention Options:} Geeignete BE-Tools
\end{itemize}

\subsection{Phase 3: Aufbereitung der Ergebnisse}

\begin{fainfobox}[Einheitliche Berichtsstruktur]
\begin{enumerate}
    \item \textbf{Ausgangslage} -- Kontext der Design-Herausforderung
    \item \textbf{Fragestellung} -- Präzise Formulierung des Problems
    \item \textbf{Verhaltensökonomische Analyse} -- Detaillierte Untersuchung
    \item \textbf{Design-Empfehlungen} -- Konkrete Handlungsoptionen
\end{enumerate}
\end{fainfobox}

\newpage
% =============================================================================
\section{Unser Leistungsangebot}
% =============================================================================

\subsection{Übersicht der Module}

\begin{center}
\begin{tabular}{p{2.5cm}p{5.5cm}p{5cm}}
\toprule
\rowcolor{fadarkblue}
\fahead{Modul} & \fahead{Eignung} & \fahead{Umfang pro Monat} \\
\midrule
\rowcolor{fawhite}
\textbf{Small} & Einzelne Design-Fragestellungen & 1 umfassend ODER 2 kurz \\
\rowcolor{falightgray}
\textbf{Medium} & Mehrere Fragestellungen & 2 umfassend ODER 3 kurz \\
\rowcolor{fawhite}
\textbf{Large} & Maximale Flexibilität & Nach Absprache \\
\bottomrule
\end{tabular}
\end{center}

\vspace{1cm}

% -----------------------------------------------------------------------------
\subsection{Modul Small}
% -----------------------------------------------------------------------------

\textbf{Ziele:}
\begin{itemize}
    \item[$\checkmark$] Klarheit schaffen bei präzisen Fragestellungen
    \item[$\checkmark$] Design-Entscheidungen absichern
    \item[$\checkmark$] Risiken frühzeitig erkennen
\end{itemize}

\textbf{Leistungen:}
\begin{itemize}
    \item 1 umfassende Analyse (8--10 Seiten) ODER 2 Kurzanalysen (3--5 Seiten) pro Monat
    \item 3--4 Wochen Bearbeitungszeit
    \item Auswertung UBS-Unterlagen, wissenschaftlich fundierte Analyse
\end{itemize}

\vspace{0.3cm}
\begin{center}
\fcolorbox{fadarkblue}{falightgray}{%
\parbox{0.85\textwidth}{%
\centering
\textbf{Dauer:} 3 Monate\\[0.3cm]
\textbf{\Large Preis: CHF 12'000.-}
}%
}
\end{center}

\newpage
% -----------------------------------------------------------------------------
\subsection{Modul Medium}
% -----------------------------------------------------------------------------

\textbf{Ziele:}
\begin{itemize}
    \item[$\checkmark$] Themen verbinden -- Zusammenhänge sichtbar machen
    \item[$\checkmark$] Prioritäten setzen -- fundierte Entscheidungsbasis
    \item[$\checkmark$] Design-Optionen vergleichen
\end{itemize}

\textbf{Leistungen:}
\begin{itemize}
    \item 2 umfassende Analysen (8--10 Seiten) ODER 3 Kurzanalysen (3--5 Seiten) pro Monat
    \item 3--4 Wochen Bearbeitungszeit
    \item Auswertung UBS-Unterlagen, wissenschaftlich fundierte Analyse
\end{itemize}

\vspace{0.3cm}
\begin{center}
\fcolorbox{fadarkblue}{falightgray}{%
\parbox{0.85\textwidth}{%
\centering
\textbf{Dauer:} 3 Monate\\[0.3cm]
\textbf{\Large Preis: CHF 18'000.-}
}%
}
\end{center}

\vspace{1cm}

% -----------------------------------------------------------------------------
\subsection{Modul Large}
% -----------------------------------------------------------------------------

\textbf{Ziele:}
\begin{itemize}
    \item[$\checkmark$] Flexibel auf kurzfristige Fragestellungen reagieren
    \item[$\checkmark$] Laufende Begleitung als Sparringspartner
    \item[$\checkmark$] Strategische Wirkung absichern
\end{itemize}

\textbf{Leistungen:}
\begin{itemize}
    \item Laufende Bearbeitung in flexibler Tiefe und Frequenz
    \item Kurzfristige ad hoc Bearbeitung möglich
    \item Periodische Executive Summaries
\end{itemize}

\vspace{0.3cm}
\begin{center}
\fcolorbox{fadarkblue}{falightgray}{%
\parbox{0.85\textwidth}{%
\centering
\textbf{Dauer:} Fortlaufend\\[0.3cm]
\textbf{\Large Preis: nach Absprache}
}%
}
\end{center}

\newpage
% =============================================================================
\subsection{Leistungsübersicht}
% =============================================================================

\begin{center}
\begin{tabular}{p{7cm}ccc}
\toprule
\rowcolor{fadarkblue}
\fahead{Leistungsbausteine} & \fahead{S} & \fahead{M} & \fahead{L} \\
\midrule
\rowcolor{fawhite}
Kick-off \& Zielklärung & $\checkmark$ & $\checkmark$ & $\checkmark$ \\
\rowcolor{falightgray}
Auswertung UBS-Unterlagen & $\checkmark$ & $\checkmark$ & $\checkmark$ \\
\rowcolor{fawhite}
Wissenschaftlich fundierte Analyse & $\checkmark$ & $\checkmark$ & $\checkmark$ \\
\rowcolor{falightgray}
Strukturierter Analysebericht & $\checkmark$ & $\checkmark$ & $\checkmark$ \\
\rowcolor{fawhite}
Standardumfang (2 kurz / 1 umfassend) & $\checkmark$ & & \\
\rowcolor{falightgray}
Erweiterter Umfang (3 kurz / 2 umfassend) & & $\checkmark$ & \\
\rowcolor{fawhite}
Ad hoc Flexibilität & & & $\checkmark$ \\
\rowcolor{falightgray}
Executive Summaries & & & $\checkmark$ \\
\rowcolor{fawhite}
Laufende Begleitung & & & $\checkmark$ \\
\midrule
\rowcolor{falightgray}
\textbf{Dauer} & \textbf{3 Mt.} & \textbf{3 Mt.} & \textbf{Fortl.} \\
\rowcolor{fadarkblue}
\fahead{Preis} & \fahead{12k} & \fahead{18k} & \fahead{n.A.} \\
\bottomrule
\end{tabular}
\end{center}

\newpage
% =============================================================================
\section{Der Nutzen für das CoE Experience Design}
% =============================================================================

\subsection{Konkrete Unterstützung im Arbeitsalltag}

\begin{itemize}
    \item[$\checkmark$] \textbf{Schnelle Orientierung:} Fundierte Antworten auf Design-Fragen
    \item[$\checkmark$] \textbf{Reduktion von Unsicherheit:} Transparente Analysen statt Best Practices
    \item[$\checkmark$] \textbf{Zeitersparnis:} Präzise Berichte statt langer Diskussionen
\end{itemize}

\subsection{Strategischer Mehrwert}

\begin{itemize}
    \item[$\checkmark$] \textbf{Positionierung stärken:} CoE XD als Quelle evidenzbasierter Empfehlungen
    \item[$\checkmark$] \textbf{Capability Building fortführen:} Methoden aus UBS\_036 kontinuierlich anwenden
    \item[$\checkmark$] \textbf{Messbare Wirkung:} Dokumentierte Analysen zeigen Wertbeitrag
\end{itemize}

\subsection{Nutzen für UBS Kund:innen}

\begin{itemize}
    \item[$\checkmark$] \textbf{Bessere Kundenerlebnisse:} Design basiert auf Verhaltensverständnis
    \item[$\checkmark$] \textbf{Weniger Friction:} Barriers werden systematisch abgebaut
    \item[$\checkmark$] \textbf{Höhere Adoption:} Features werden tatsächlich genutzt
\end{itemize}

\begin{fasuccessbox}[Zusammenfassung]
UBS Switzerland gewinnt mit diesem Angebot ein Werkzeug, das den Arbeitsalltag des CoE Experience Design erleichtert, die strategische Positionierung des Teams stärkt und evidenzbasiertes Design als Standard etabliert.
\end{fasuccessbox}

\newpage
% =============================================================================
\appendix
\section{Verbindung zum Capability Building (UBS\_036)}
% =============================================================================

\begin{center}
\begin{tabular}{p{6cm}p{6cm}}
\toprule
\rowcolor{fadarkblue}
\fahead{UBS\_036 (abgeschlossen)} & \fahead{Frage-Abo (neu)} \\
\midrule
\rowcolor{fawhite}
Grundlagen vermittelt & Anwendung im Alltag \\
\rowcolor{falightgray}
BE-Sprache etabliert & BE-Analysen durchgeführt \\
\rowcolor{fawhite}
Ambassadors ausgebildet & Ambassadors unterstützt \\
\rowcolor{falightgray}
Experimentation Mindset & Experimente begleitet \\
\bottomrule
\end{tabular}
\end{center}

\section{Beispiele für typische Fragestellungen}

\textbf{Journey-Design:}
\begin{itemize}
    \item «Wo in der Onboarding-Journey verlieren wir die meisten Neukund:innen?»
    \item «Wie sollte der Wechsel von App zu Filiale gestaltet sein?»
\end{itemize}

\textbf{Touchpoint-Optimierung:}
\begin{itemize}
    \item «Welche Defaults sollten wir im Anlageberater setzen?»
    \item «Wie reduzieren wir die Komplexität bei der Produktauswahl?»
\end{itemize}

\textbf{Feature Adoption:}
\begin{itemize}
    \item «Warum wird Feature X kaum genutzt?»
    \item «Wie steigern wir die Nutzung der Budgeting-Funktion?»
\end{itemize}

\section{Über FehrAdvice \& Partners}

FehrAdvice \& Partners ist das führende Beratungsunternehmen für Behavioral Economics im deutschsprachigen Raum. Gegründet von Prof. Ernst Fehr (Universität Zürich).

\begin{itemize}
    \item 15+ Jahre Erfahrung in Behavioral Economics Beratung
    \item 200+ Projekte in verschiedenen Branchen
    \item Wissenschaftliches Fundament durch Prof. Ernst Fehr
\end{itemize}

\vfill

\begin{center}
\textcolor{fadarkblue}{\rule{0.8\textwidth}{1pt}}\\[0.5cm]
\textbf{FehrAdvice \& Partners AG}\\
Klausstrasse 4 | 8008 Zürich\\
www.fehradvice.com
\end{center}

\end{document}
